% Deutsche Parallelfassung zu foundations/part_intro.tex

\chapter*{Mathematische Grundlagen}
\phantomsection
\addcontentsline{toc}{chapter}{Mathematische Grundlagen}

Die Krümmungsänderung hat die Argumentation an eine genaue Schwelle geführt.
Die Thesis muss nun beschreiben, was gleich bleibt, was sich ändert, welche
Größen daraus folgen und unter welchen Qualifikationen eine Bewertung gültig
ist. Unterscheidungen in natürlicher Sprache sind notwendig, reichen jedoch
für Ableitung, Vergleich oder Implementierung nicht aus.

Dieser Teil gibt daher ausgewählten ingenieurtechnischen Abhängigkeiten eine
mathematische Form. Er beginnt mit Informationsmodellierung und begrifflicher
Struktur, weil eine Formel nur nützlich ist, wenn die Rollen ihrer Ein- und
Ausgaben klar sind. Anschließend formuliert er die maßgeblichen Axiome und die
Notation, entwickelt Zustands- und Operatorbeziehungen und erklärt, wie
intrinsische Konstruktion durch metrische Realisierung messbare räumliche
Bedeutung erhält.

Mathematik entscheidet weder über die Identität eines Engineering Objects,
noch nimmt sie Wissen an oder übt Ingenieurautorität aus. Sie kann
Abhängigkeiten ausdrücklich machen, Zustände ableiten, eine
Krümmungsänderung fortpflanzen, eine Beobachtung mit einem Ziel vergleichen
und eine Nebenbedingung bewerten. Die in Teil~I eingeführten
Unterscheidungen müssen jede dieser Operationen überstehen.

Die orientierende Frage lautet:

\begin{quote}
Welche formalen Abhängigkeiten machen Folgen berechenbar, ohne das Alignment,
seine Zustände, seine Repräsentationen, die Evidenz darüber und die
möglicherweise folgende Entscheidung miteinander zu verwechseln?
\end{quote}

Das erste Kapitel überschreitet diese Schwelle anhand der bereits aus der
Praxis bekannten Repräsentationen.
